\documentclass[USenglish,twocolumn]{article}

\usepackage[utf8]{inputenc}%(only for the pdftex engine)
%\RequirePackage[no-math]{fontspec}%(only for the luatex or the xetex engine)
\usepackage[big]{dgruyter}
%\graphicspath{{figures/}} % path to folder with figures
\usepackage{xcolor}
\usepackage{float}
\usepackage{lipsum} 
%\usepackage{blindtext}
%\usepackage[showframe]{geometry}
\usepackage{graphicx}
\usepackage[super]{nth}
\usepackage{gensymb}
\usepackage{soul}

\begin{document}
 
\articletype{Research Article{\hfill}Open Access}
\author[1]{Per Lindh}
\author*[2]{Polina Lemenkova}
  \affil[1]{$^1$ Swedish Transport Administration, 
Department of Investments, Technology and Environment, Neptunigatan 52, Box 366, SE-201-23, Malmö, Sweden \\ 
	$^2$ Lund University, Lunds Tekniska Högskola LTH (Faculty of Engineering), Department of Building and Environmental Technology, Division of Building Materials, Box 118, SE-221-00, Lund, Sweden, E-mail: per.lindh@byggtek.lth.se}
  \affil[2]{$^3$ Université Libre de Bruxelles (ULB), École polytechnique de Bruxelles (Brussels Faculty of Engineering), Laboratory of Image Synthesis and Analysis (LISA). Campus de Solbosch -- CP 165/57, Av. F. D. Roosevelt 50, B-1050 Brussels, Belgium, E-mail: polina.lemenkova@ulb.be}

  \title{\huge Seismic monitoring of strength in stabilized foundations by P-wave reflection and downhole geophysical logging for drill borehole core}
  
 \runningtitle{Seismic monitoring of strength by P-wave reflection and geophysical logging}
 \runningauthor{P. Lindh, P. Lemenkova}

\begin{abstract}
{Evaluating the subground properties during the initial stage of a construction of building is important in order to estimate the suitability of soil quality to the technical requirements of bearing capacity, resistance to stress and strength. This study presented the evaluation of the geotechnical properties of soil intended for construction of Max IV facility of Lund University, performed in fieldwork and laboratory. The \emph{in-situ} methods included drilling boreholes, core sampling and assessment, crosshole measurements and borehole logging. The laboratory measurements were performed at Swedish Geotechnical Institute (SGI) and combined seismic measurements of drill cores, determination of the Uniaxial Compressive Strength (UCS) and examination of material property: sieve analysis and natural moisture content. The UCS was evaluated with regard to velocities of elastic P-waves. The synchronous light test by X-ray diffraction was performed for qualitative analysis of mineral composition of samples. The study applied integrated approach of the diverse geophysical methods to solve practical tasks on evaluation of foundation strength and geotechnical parameters. This study demonstrated the benefits of integrated seismic and geophysical methods applied to soil exploration in civil engineering for testing quality of foundation materials.}
\end{abstract}
  \keywords{seismicity; elastic wave; X-ray diffraction}
 
% seismicity, elastic wave, X-ray diffraction, mineral composition, gamma-ray spectrometry log, borehole, drilling, crosshole, core sampling
%  \classification[PACS]{}
 % \communicated{...}
 % \dedication{...}

  \journalname{Journal of the Mechanical Behavior of Materials}

\DOI{DOI}
  \startpage{1}
  \received{Aug 25, 2022}
  \revised{Mar 24 2023}
  \accepted{..}

  \journalyear{2023}
  \journalvolume{x}
  \journalissue{x}
  
\maketitle

\section{Introduction}

Geophysical methods offer an effective tool for evaluation of physicochemical and mechanical parameters of soil in civil engineering. Information regarding the lithological structure and mineral content, geochemical and geological properties of soil can be extracted from the core drilling and geophysical borehole logs \cite{MAUCEC2009194,JING2021101534}. In recent years, research shows that geophysical methods of data processing are certainly required in cases when non-destructive methods are advisable. For instance, parameters of acoustic velocities of seismic waves penetrated through the soil specimens can be used as a dataset for estimation of soil strength in a laboratory using the reflectivity of P-waves \cite{Lindh202102,Lindh202210}. Further, data processing aimed to investigate strength, elasticity and mineral properties of soil specimens can be performed using the acoustic imaging and seismic testing by measuring the velocities of P- and S-waves \cite{Tang2009,Tronicke,Hadden} as well as X-ray-diffractions for testing the mineral content of soil \cite{Wang}. 

Evaluation of strength properties of soil is an essential requirement for construction works for safety and stability of buildings. Prior to the constructions and engineering works, clayey soil should be reinforced using binders. In the environmental and climate conditions of Sweden, weak expansive clayey soil is subject to seasonal swelling and shrinking. This can lead to a serious damage of civil engineering objects, roads and infrastructure. To avoid this, clayey soil should be stabilized using binders prior to its use in engineering works. Stabilization may also include the environmental aspects, such as cleaning the material from chemicals, pesticides \cite{CONTESSI2021141826,Lemenkova202222} and toxic hazardous materials  \cite{Lindh202201}. 

Stabilization of soil by binders results in the increased strength, which can be evaluated using the uniaxial compressive strength (UCS) tests. Binders usually includes such materials as cement \cite{Lindh2004,HORPIBULSUK20102011}, lime \cite{OKYAY201034,JAUBERTHIE2010395}, slag \cite{Lindh2021a}, cementitious materials \cite{YOON20102322} in various proportions and content \cite{Lindh2001}. Additional methods of evaluation of soil quality and parameters include the geophysical data processing and interpretation of the borehole logs \cite{LIU2022107048,Rydenetal2016,COMINA2021106278}. Advances in the field of applied geophysics and data analysis made an important application become real: evaluating the soil characteristics through quantitative estimation of its parameters using the nondestructive geophysical methods. Such approaches include, for instance, the inverse geophysical methods which are based on the correlation of density and porosity of core with the response from the elastic seismic waves. The evaluation of these data can be performed using an integrated approach of interpretation of the crosshole measurements with acoustic testing of the core specimens. The aim of such applications is to evaluate the soil strength using measured seismic velocities of P-waves of the drill core. 

Many civil engineering tasks can benefit from the use of applied geophysical methods. For instance, the interpretation of borehole logging enables to obtain information on structural characteristics and mineral composition of soil which can be determined from data analysis. The borehole logging presents an effective tool for geotechnical exploration for ground measurements inside the boreholes, widely used for soil analysis and reported in a variety of studies \cite{KALERIS2015135,DIAS2020107099,YU2022127817,YU2021126010}. Probe sensors of the borehole logging system measure the physical properties of the ground such as density, and porosity, magnetic and radioactivity properties, electrical resistivity, seismic velocity, 
\cite{KILLEEN2015447}. These data are useful for evaluation of the soil quality and its suitability for the construction works.

The X-ray diffraction (XRD) is another advanced method of evaluating the soil structure. This method is especially effective for the analysis of soil structure, including the evaluation of microfabrics, orientation of crystals and identification of mineral content \cite{Denaix,Das,Righi}. The XRD applies the interference of the monochromatic X-rays in a crystalline specimen. The crystalline structure of soil diffracts the X-rays and responses to an X-ray diffractometer by a generated interference image. This pattern contains the information regarding the molecular structure of soil \cite{ALDERTON2021520}. The benefits of the XRD for soil analysis is that it is a straightforward and functional method for qualitative identifying of the soil structure, for quantifying the relative proportions of the minerals in a soil sample and which can be used for detecting the soil evolution phases \cite{ADAMS2019391}.

Optical images of the borehole wall of the drilling core in artificial colours are produced by the core scanner devices which allow to record and visualise the wall sections of drilling cores. The acoustic televiewer is an essential feature of the borehole imaging based on the ultrasonic pulse-echo configuration \cite{paillet1990borehole}. Its approach consists in plotting the borehole by the ultrasound, which is transmitted from the piezoelectric transducer in the borehole. The waves are reflected on the borehole by the mirror and measured as a log by the receiver, which records the amplitudes and the travel time of the signals \cite{RONQUILLOJARILLO20181}. These data record information regarding the fracture strike and the dip \cite{Paillet}. The acoustic data recorded by the borehole imaging system are used as an indication of the soil properties: character, relation, and orientation of the lithologic and structural features of soil specimen \cite{WILLIAMS2004151}. 

New approaches in the acoustic borehole imaging are constantly developing with aim to achieve the optimum performance of the device: to improve speed and precision of the data processing workflow \cite{schepers1991development,DELTOMBE2004161}, to increase the quality of imaging \cite{ALSIT2015139}, to enhance the recognition of echoes and separate noise \cite{WEDGE2017157}, to adjust to the real-time borehole conditions \cite{Jerram}. In soil analysis, digital image and data processing provides an important source of information which can be used as to control and improve the interpretation of geologic measurements obtained from borehole drilling \cite{Kaellenetal2014,Kallenetal2016,Lemenkov2021c}.

The geologic impulse response from gamma ($\gamma$) rays in various beds of a borehole enables to understand the geological nature and geophysical properties in complex sequences of the drilled layers \cite{Conaway}. Gamma-ray logging of the boreholes and spectrometer system have been reported for the determination of mineral compositions in a drill core \cite{Ajayi} and concentrations of metals and radioactive elements ($^{238}U$, $^{133}Ba$, $^{232}Th$, $Mn$) \cite{Loevborg,Mikesell}. For instance, gamma-ray logs enable to perform the granulometric estimation of the clay content in a soil sample, due to the high porosity and low resistivity of clay, which, together with its mineralogical composition (potassium) causes high gamma-ray emissivity \cite{Bhuyan,DIAZCURIEL2021104481}. The compensated gamma-gamma ($\gamma-\gamma$) type density is a log corrected for the negative effects of the gamma rays naturally occurring in radioactive elements ($^{238}U$, $^{232}Th$), which is used for measuring the density of the borehole core layers \cite{osti_5961810,osti_5614551}.

The information regarding the geotechnical properties of soil intended to foundation constructions is normally acquired by the inspection of strength (UCS) \cite{LindhWinter,Lemenkov2021b,Lemenkov2021a}. Strength of soil also provides data on the quality of stabilization works. The use of the compression machine for UCS tests is usually applied to detect precise information of soil behaviour in strain-stress conditions. More advanced geophysical methods include the non-destructive approaches using the evaluation of the velocity of seismic P-waves passed through soil specimens for analysis of strength. 

The integration of the geophysical and engineering geological methods is important, since it enables to use the benefits of both approaches and apply the diverse tools for a complex assessment of soil properties prior to construction works. For instance, the results on seismic tests show the correlation between the UCS and P-wave velocities which is useful to control the quality of soil stabilization, while borehole logging provides the information regarding the local properties of the tested soil specimens intended for a construction of foundation of buildings.

This study presents the results of work aimed at testing and improving the parameters of soil prepared for construction of foundation of the large building facility. These include the geophysical borehole logging for the core analysis, tests on soil stabilization and geophysical investigations. The requirements for soil included high strength, low sensitivity to vibrations, high bulk density and low porosity. The plan of this article is as follows. The scope of this study and \emph{in-situ} condition of a real project were presented in Section \ref{sec:2}. Section \ref{sec:3} reports the workflow of fieldwork, presents selected examples of the core sampling and borehole logging. Section \ref{sec:4} defines the laboratory tests, including the UCS and seismic tests. Section \ref{sec:5} presents the details of data processing and describes the obtained results. Finally, the discussion of the results is presented in Section \ref{sec:6}, where the comments on the performance of the methods are included. Some remarks concerning the future works are provided in the last Section \ref{sec:7}.

\section{Background and purpose}\label{sec:2}
This paper aims at testing soil strength in a planned construction of building using applied geophysical methods and engineering geological methods. The test area belongs to the synchrotron radiation light laboratory Max IV (\href{https://www.maxiv.lu.se/}{www.maxiv.lu.se/}) at Brunnsh\"og in Lund, coordinated by Lund University, Fig. \ref{Fig:01}. 

\begin{figure}[H]
 \centering
    \includegraphics[width=0.50\textwidth]{Fig_01.jpg}\\
  \caption{Fieldwork measurements at Max IV. Photo: Per Lindh.}\label{Fig:01}
\end{figure}

An essential parameter for the performance of Max IV facility is related to the vibration characteristics because soil foundation should be resisted to the high level of vibrations. To evaluate the quality of foundation that should meet the technical requirements regarding the vibration, a series of simulations were carried out using geophysical methods. The purpose of the project is to improve the quality of the subsoil prepared for construction works and to test the results of stabilization through the evaluation of the compressive strength characteristics of the stabilized soil. The investigation of the properties of the entire stabilized layer enables to evaluate how well the selected production model succeeded in creating a homogeneous monolith of the stabilized soil material using core drilling, Fig. \ref{Fig:02}. 

\begin{figure}[H]
 \centering
    \includegraphics[width=0.45\textwidth]{Fig_02.jpg}\\
  \caption{Core drilling. Photo: Per Lindh.}\label{Fig:02}
\end{figure}

The drilling of the borehole core was performed to select soil samples, stabilized and investigated by the downhole geophysical logging and seismic methods. This facility which has a unique basement foundation of the building consisting of $160,000 m^3$ of clay moraine which required to be stabilized and tested for strength. 

The facility includes the independent experimental stations, which operate with a wide range of the experimental techniques, including the crystallography, electron spectroscopy, nanolithography and photo-nuclear experiments. One of the prerequisites for achieving the stable performance of the foundation of Max IV is to reduce the effects from vibrations on soil. Therefore, this work consisted of reducing the plant's sensitivity of foundation to vibrations by increasing its rigidity, so that the components in accelerators and experimental stations could move in time. The foundation also has a certain reducing effect on vibrations from the internal and external sources. The required condition for vibration level is to be below 20-30 nm RMS displacement for frequency >5Hz. 

The research included the following objective: 

\begin{itemize}
	\item to achieve the homogeneity of the subsoil prepared for basement construction;
	\item to evaluate the strength in the bonds in boundaries between the different layers of soil;
	\item to estimate the difference in density between the surface and deep-lying material through additional compaction provided by the overlying clay;
	\item to evaluate how the soil properties were improved in the stabilized specimens after curing, compared to the initial non-stabilized raw soil samples;
	\item to investigate the binding products on the stabilized material using \emph{in-situ} measurements in Max IV;
	\item to analyse the differences between the seismic results from the cross hole and seismic tests on the drill cores.
\end{itemize}

As a result of the performed works, the bearing capacity of soil was improved which was proved by a series of tests, as presented and described below in sections below.

%\newpage
\section{Fieldwork measurements}\label{sec:3}
The measurements were performed in connection with the construction works aimed to improve the quality of soil through stabilization and to evaluate the results of soil stabilization through a series of seismic tests. 

\subsection{Core sampling}

The field tests started with three core samples being drilled as a cylindrical sections through the stabilized layer of soil down to the underlying naturally deposited clay moraine. The drilling was performed using the S Geobor, a wireline triple tube core barrel drilling system developed by Terracore, \cite{Epiroc}. The hole bit diameter is 146 mm; the core inner diameter is 102 mm. Fig. \ref{Fig:02}. The drilling was carried out to a depth of 5 m, i.e., an under-drilling was of ca. 1 m below the stabilized layer. No core losses were reported in any of the three core drill holes. The examples of the drill cores are shown in Fig. \ref{Fig:03} and Fig. \ref{Fig:04}.

\begin{figure}[H]
 \centering
    \includegraphics[width=0.45\textwidth]{Fig_03.jpg}\\
  \caption{Drill core from the borehole no. 2, depth: between 4--5 m. The right part of the drill core consists of natural clay moraine, which is more plastic than the stabilized clay moraine. Photo Per Lindh.}\label{Fig:03}
\end{figure}

\begin{figure}[H]
 \centering
    \includegraphics[width=0.45\textwidth]{Fig_04.jpg}\\
  \caption{Drill core from the borehole no. 2, depth: between 1--2.5 m. The drill core shows homogeneous stabilization along the entire stretch. Photo: Per Lindh.}\label{Fig:04}
\end{figure}

\subsection{Stabilization}
Stabilization of soil is a method that works best for improving its resistance towards the internal and external sources of vibration and load. Natural soil was stabilized beneath the two storage rings with the hydrated lime and ground granulated blast furnace slag (GGBFS) down to a depth of 4 m. The stabilization was carried out in layers of 3 m with an undercut in the underlying layer between 3-5 cm to ensure a homogeneous matrix. The continuous quality control was performed of each stabilized soil layer. 

\subsection{Crosshole measurements}
The tests on the crosshole measurements were based on the three core samples collected from the ground surface down to the bottom of the stabilized layer (to 3.15 m). The specimens of the drill cores were mapped and photographed, after which selected samples were measured for P-wave velocity, compressive strength, density and deformation. The specimens were tested at the diffraction beam tube BLI 711 at MAX-lab, Öle Römers väg, Lund. The crosshole measurements were performed with a new type of equipment which has extended possibilities to measure shear wave and compression wave velocity in the material between the boreholes. The system consists of a transmitter unit that generates shear or compression waves using two receiver units (geophones), Fig. \ref{Fig:05}. 

\begin{figure}[H]
 \centering
    \includegraphics[width=0.45\textwidth]{Fig_05.jpg}\\
  \caption{Schematic sketch of crosshole measurements on Max IV.}\label{Fig:05}
\end{figure}

The tests were performed in accordance with standard ASTM D4428/D4428M-14, \cite{ASTM2016}. The crosshole measurements show the achieved requirements for the soil quality in terms of velocities of horizontally travelling seismic waves in soil: primary compression wave (P-wave) and secondary shear wave (S-wave). Since stabilized soil was considered sufficiently strong, no reinforcement of the borehole walls was carried out except by a thin layer filled on top of the stabilized layer, Fig. \ref{Fig:06}. 

\begin{figure}[H]
 \centering
    \includegraphics[width=0.45\textwidth]{Fig_06.jpg}\\
  \caption{Crosshole measurements on Max IV. The seismic source was located in the borehole No. 1 on the far right of the image. Plastic pipes go down to the stabilized layer. Photo: Per Lindh.}\label{Fig:06}
\end{figure}

The design criterion was a minimum shear wave speed ($W_s$) of 900 m/s. Fig. \ref{Fig:07} shows a $W_s$ between 1000--1350 m/s (\nth{3} curve from the left, red). The coefficient of lateral expansion is shown in the Poisson's ratio (\nth{5} curve left, red), with the lowest values (0.14) on the horizon at 1.20 m.

%\newpage
\subsection{Borehole logging}
The geophysical borehole logging has been performed using the technical tools of Ramboll Group AS. The examples of the detailed and precise analysis of the variations in physical parameters of core are presented in Fig. \ref{Fig:08}, Fig. \ref{Fig:09} and Fig. \ref{Fig:10}. The logging has been carried out in the core boreholes with aim to derive lithological, structural and rock mechanical information using geophysical methods. Geophysical borehole logs in Fig. \ref{Fig:08}, Fig. \ref{Fig:09} and Fig. \ref{Fig:10} also show the data for lithology analysis by a combined evaluation of the log responses and acoustic measurements. 

The borehole logging system collected information on core using probes with embedded sensors. The data captured by the borehole logging system were transmitted to the computer using a logging cable. The system included a GPS and a depth encoder for the 3D coordinate capture. The results of the data processing include the borehole logging which shows several plots of sensor outputs visualised in a horizontal coordinate system versus depth on the vertical axis, as showed in Fig. \ref{Fig:08}, Fig. \ref{Fig:09} and Fig. \ref{Fig:10}.

Using borehole logging has important benefits for soil testing: 
\begin{enumerate}
	\item First, the reflected impulse of waves is directly related to the elastic and material properties of soil. This provides data on soil structure which can be received both using the direct measurements and indirect evaluation. The latter includes solving the inverse problem to estimate hidden parameters of soil from its related properties (e.g., porosity-density links). 
	\item Second, the acoustic testing device of the borehole logging is highly sensitive to the microstructure of the ground. Thus, the recording of logging can be used for recognition of the defects, detection of fractures and cracks in layers prepared for foundation construction.
\end{enumerate}

\begin{figure*}
 \centering
    \includegraphics[width=0.95\textwidth]{Fig_07.jpg}
  \caption{Results from the crosshole measurements.}\label{Fig:07}
\end{figure*}

The information is confined to a limited distance of 4.95 m tested around the boreholes drilled for Max IV. The structural analysis of core derived from borehole logging comprised data on several following parameters illustrated as vertical cross sections: 

\begin{enumerate}
	\item depth from the ground level, 
	\item natural $\gamma$ rays, 
	\item inclination from the vertical level, 
	\item caliper probe 1-arm,
	\item compensated $\gamma-\gamma$ density,
	\item porosity calculated from density,
	\item filtered radius image from the borehole wall (360\degree),
	\item 3D caliper with borehole acoustical image,
	\item amplitude (360\degree) of acoustical image of borehole wall,
	\item hardness computed from the acoustical median of the borehole wall.
\end{enumerate}

The inclination shows the angle of vertical inclination formed between the geodetic point of measurement (located device) and the exact coordinate location of the core. The inclination phenomena is owing to the physics of the Earth where the magnetic lines that are not parallel to the surface. This results in a slightly titled position of the compass with inclined needle. This information is plotted on the graphs in azimuth \degree. The caliper probe 1-arm is a mechanical device, used to measure the internal diameter of a borehole and to press a section of a probe against the side of the borehole for density logging.
	
The results of the analysis and characterisation of soil structures are shown in Fig. \ref{Fig:08}, Fig. \ref{Fig:09} and Fig. \ref{Fig:10}, which give the information regarding the natural gamma, calliper, acoustic televiewers, density and porosity, which have been defined by a combined interpretation of the core images and acoustic images of the borehole. The hardness of the borehole wall was calculated using the results of the acoustic measurements. The results of the measurements of boreholes 2 and 3 show a variation in values at the depth of 2.5 m, while for borehole 2 it is also noted at 3.3 m. This variation in caliper results is also seen as a change in the hardness of the borehole wall. 

During the examination of the drill cores, a thin layer (ca. 10 mm) of the not stabilized soil was detected which  well corresponds with the 2.5 m level of depth. The most likely reason for such variation is that the milling down in the underlying layer was too small and therefore this layer has a lower strength. This difference in strength is not visible in other three boreholes and hence probably has a small extent for a selected borehole. The difference in soil strength does not affect the function of the foundation at the frequencies where the construction is sensitive to vibrations since these wavelengths are much longer (bigger) than this layer.

\begin{figure*}
 \centering
    \includegraphics[width=0.95\textwidth]{Fig_08.jpg}
  \caption{Results from the borehole logging for borehole no. 1.}\label{Fig:08}
\end{figure*}

\begin{figure*}
 \centering
    \includegraphics[width=0.95\textwidth]{Fig_09.jpg}
  \caption{Results from the borehole logging for borehole no. 2.}\label{Fig:09}
\end{figure*}

\begin{figure*}
 \centering
    \includegraphics[width=0.95\textwidth]{Fig_10.jpg}
  \caption{Results from the borehole logging for borehole no. 3.}\label{Fig:10}
\end{figure*}

\section{Laboratory tests}\label{sec:4}
Laboratory investigations were performed in the Swedish Geotechnical Institute (SGI) and included (i) soil analysis (grain distribution and natural water content), (ii) seismic measurements and (iii) UCS strength tests of the drill cores (Fig. \ref{Fig:11}). These investigations were based on the three core samples that were carried out from the ground surface down to the bottom of the stabilized layer. The drill cores were mapped and photographed, after which the samples were selected for further determination of the P-wave velocity, compressive strength and density. Furthermore, the material samples were tested at the diffraction beam tube BLI 711 at MAX-lab, Öle Römers väg, Lund.

\subsection{Compressive strength of drill cores}

%\textcolor{red}{\lipsum[1]}
The soil parameters were improved through stabilization, which was performed using binders adjusted to the specific properties of soil: GGBFS and CaO, as also tested in previous studies showing the experiments on various binder combinations on the results of soil properties \cite{Lindh202232,Xu}. Upon stabilization, the compressive strength of soil was evaluated. To ensure a good performance of foundation planned to be constructed on the given soil, the advanced geophysical methods were applied. These included drilling, core sampling and assessment, crosshole measurements and borehole logging for testing soil. These data were then processed in the laboratory of SGI. The latter included seismic tests of drill cores, UCS and examination of material property of soil through sieve analysis for grain distribution and natural moisture content.

The UCS was tested according to the standards of the deformation controlling test by the mechanical press MTS, Fig. \ref{Fig:11}. The deformation rate was set constant to 2 mm/min, as tested specimens had a diameter of ca. 100 mm and a height of ca. 200 mm. Fig. \ref{Fig:11} shows a photo of the mechanical press MTS used for testing the UCS of the specimen collected in the Bore 3:3 at depth: 2.95--3.15 m. The samples were compressed in a 500 kN. The 1\% per minute was in the axial direction.

Since stabilized soil hardens differently depending on the stress level, this process could not be controlled in the \emph{in situ} measurements which showed a much slower curing process. The difference in hardening of soil in the laboratory and the \emph{in situ} conditions can be partly explained by the difference in the temperature between the outdoor air conditions and the field setting. However, the differences in the overlay pressure might also add their impact, which could contribute to the strength development for the specimens collected in field. In fact, the test specimens were not stored with a successively increased ambient pressure as in the real design which might have the impact on the results of soil hardening.

\begin{figure}[H]
 \centering
    \includegraphics[width=0.40\textwidth]{Fig_11.jpg}\\
  \caption{The mechanical press MTS used for UCS test. The axial resonant frequency -- 4516 Hz; flex mode resonant frequency -- 2902 Hz. Photo: Per Lindh.}\label{Fig:11}
\end{figure}

Other parameters aimed at evaluation of the stiffness of soil which is being assessed using Young's modulus (E) and rigidity modulus (G). The Young's modulus evaluates the elastic properties of soil in tension or compression as a ratio between the tensile or compressive stress and the axial strain. In turn, the rigidity evaluates the shear stress of soil and is denoted by the G modulus which depends on the degree of distortion of basement during deformation that results in shear stress being changed to the shear strain. These parameters are improved during the process of stabilization and solidification of weak clayey soil.

\subsection{Sieve analysis and moisture content}
The deformation performance of soil as a response to stress is a complex and non-linear parameter, controlled by many factors. Mineralogy plays an essential role in strength and stress-strain behaviour of soil. It includes the porosity and ratio of voids that correlate with soil stiffness; plasticity and rheology that affect the mechanical properties of soil; structure and fabric related to the interparticle contact in soil. Since inner properties of soil influence its response to stress and strain, the analysis of soil structure is a fundamental part of the construction. 

The sieve analysis test was performed according to the standard ASTM D7928-21E1 \cite{D7928} for particle size distribution of fine-grained clayey soils by sedimentation analysis using hydrometer. The sieve analysis determined quantitatively the distribution of particle sizes of the fine-grained portions of soil samples. The fine-grained fraction of a soil was detected and discriminated within a range of other particle sizes as percentage of fraction, and the gradation curve was drawn to measure the ratio of each size of grain that is contained within the tested soil specimens. 

\subsection{Seismic \textcolor{blue}{tests} of \textcolor{blue}{the} drill cores}
Evaluating the correlations between the seismic response of the elastic waves and strength properties of the stabilized soil is a key approach in non-destructive measurements of civil engineering. Given soil specimens are of identical type, the task is to identify the level of the resonant frequency in the ultrasonic waves, such that the values of P-wave velocity correspond to the level of strength on each measured sample of soil specimen. Since soil is used as a basis for the building, this sets up the prerequisites to achieve a high performance of the ground foundation: stability, strength and reduced effects from the vibrations. 

Seismic measurements were performed as a free-free resonant column measurements of the P-wave velocity. In this test, a cylindrical sample was placed atop of a soft foam layer to simulate free boundary environment. The vibrations of P-waves were generated by a small hammer on one end of the tested specimen and measured velocities by the accelerometer which was attached to its another end. The fundamental frequency of seismic waves was measured with a ceramic shear ICP® accelerometer. The vibrational response of the sample correlates with the inner properties of its structure and thus enables to estimate strength and small-strain moduli of soil. 

This is possible due to the existing correlation between the P-wave velocities with strength of the porous material. Thus, the more dense and less porous the soil is, the higher is the velocity of propagation of the P-waves penetrating through the soil. Vice versa, the increasing porosity in soil results in the lower speed of P-waves. Such fundamental properties of the values of seismic velocity enables to use them as a robust indicator of the engineering properties of soil, and to evaluate such characteristics of soil as compressibility, permeability, elasticity and cohesion. In turn, these parameters of soil are crucial factors for evaluating its suitability as a basement prior to building construction as a reliable data on soil quality and safety of future building.

Measurements of the flex mode in tested specimens were carried out since flex mode gives a higher determination of the shear wave velocity in soil sample as it is more affected by the boundary conditions of the specimens. Comparing longitudinal resonant frequency (P-waves) and bending resonant frequency (S-waves) was performed for a quality control regarding the requirements for soil resistance to cracks and inhomogeneities in a specimen. For tested soil specimens, the relationship between the P and S waves is shown in Fig. \ref{Fig:15}.

\section{Results}\label{sec:5}

\subsection{Sieve analysis and natural moisture water content}

The results of the sieve analysis are presented in grain size distribution curve of the natural clay moraine which varied in different samples of the selected specimens, as presented in Fig \ref{Fig:12}.

\begin{figure}[H]
 \centering
    \includegraphics[width=0.50\textwidth]{Fig_12.jpg}\\
  \caption{Grain distribution of soil performed to select binders for stabilization aimed to strengthen the foundation on Max IV}\label{Fig:12}
\end{figure}

The plots of the grain distribution curve were used to analyse the type of soil, which was detected to be moraine clays. The clay content varied between 22 to 45\% of the total soil weight. Therefore, the binder recipe for stabilization was selected according to this type of soil and consisted of 50\% quicklime (CaO) and 50\% of GGBFS. This recipe of binders proved to be robust enough to handle the variations in grain distribution and effectively stabilize clayey soil. 

Natural moisture content was tested in order to obtain the information regarding the soil properties (Fig. \ref{Fig:13}), to select a proper binder for stabilization and fine an optimal a recipe to handle the variations in a clay content. The natural water ratio ($W_N$) of soil specimens varied between 11 to 23\% with an average value of ca. 18\%, Fig. \ref{Fig:13}. The experiments on the variations in water ratio within soil specimens were performed by adding extra water and evaluating the compactness of soil using the Moisture Condition Value (MCV) testing. The MCV test was used for assessing and testing the suitability of soil based on specimens being compacted by the soil settlement. The aim of this test was to evaluate the parameters of soil specimens in relation to the upper limits of moisture content.

\begin{figure}[H]
 \centering
    \includegraphics[width=0.45\textwidth]{Fig_13.jpg}\\
  \caption{Natural moisture water content ($W_N$) in the samples.}\label{Fig:13}
\end{figure}

\subsection{Correlation between the seismic wave velocity and compressive strength}

The results of the soil stabilization have been assessed using seismic velocities of the elastic waves using free-free resonance method, similar to the existing studies \cite{IsmailRyden} and reported in a series of plots illustrating the correlation between the gain in the compressive strength of soil and the wave speed. Using seismic testing is an important approach as it focuses on the advanced, non-destructive and robust evaluation of strength parameters of soil compared to the UCS traditional methods. Complicated properties of fine-grained moraine soil collected in real in-situ conditions requires that precise approaches are used for analysis of suitability of soil used as a basis for a constructed building with highly demanding requirements regarding resistivity towards vibrations, density and stiffness.

The P-wave and S-wave velocity data were integrated for the drill cores for optimisation of the strength analysis. Eight samples were tested for the analysis of UCS and measurement of P- and S-wave velocity. The requirement for test specimens was a P-wave velocity of at least 1430 m/s, which was achieved during the performed series of the experiments. The analysis of the seismic velocities of the P-wave and S-waves as well as their reflectivity measured in the boreholes demonstrated to be the essential methods for the evaluation of the strength of soil prior to the construction works.

Fig. \ref{Fig:14} shows the P-wave velocity of the soil specimens stored at 20°C which achieved the required value already after 500 hours of soil curing. Seismic measurements were carried out both on the compacted specimens collected \emph{in situ} in the field and on the surface in laboratory conditions. The resonance frequencies of the P-waves of the specimens were measured before the UCS tests in the laboratory, which demonstrated the overall increase in the velocity of waves passing through tested specimens with curing time. This well correlates with the results of similar studies that measure P-wave velocity in UCS tests for evaluation of gain in strength of soil specimens 
\cite{Aldaood,Lindh202232,Azimian}. 

\begin{figure}[H]
 \centering
    \includegraphics[width=0.50\textwidth]{Fig_14.jpg}\\
  \caption{P-wave (compression) velocity against the curing time for laboratory tested soil samples stored at 20°C.}\label{Fig:14}
\end{figure}

Since the selected adhesive has a slow-setting nature, only the initial curing phase was measured \emph{in situ}. The curing process of soil specimens was performed on the laboratory-treated specimens that were stored in a workplace of the SGI. The comparison between the core sampling and cross hole measurements shown consistent results. The stabilized soil well meets the set up requirements for a minimum of shear wave speed at least 900 m/s. Moreover, the study shown that the performed series of works on soil analysis, treatment and improvement produced a homogeneous monolith without any major deviations in material properties, which was accepted for engineering constructions. The results of the \emph{in situ} measurements demonstrate a correlation between the P- and S-values, see Fig. \ref{Fig:15}.  

\begin{figure}[H]
 \centering
    \includegraphics[width=0.50\textwidth]{Fig_15.jpg}\\
  \caption{P-wave versus S-wave velocity for drill cores}\label{Fig:15}
\end{figure}

The maximum compressive strength amounted to 2430 kPa at a deformation of about 1.75\%, as demonstrated in Fig \ref{Fig:16}. 

\begin{figure}[H]
 \centering
    \includegraphics[width=0.50\textwidth]{Fig_16.jpg}\\
  \caption{UCS vs deformation; sample bore 3:3, depth:2.95--3.15m.}\label{Fig:16}
\end{figure}

The majority of the tested specimens were selected from a depth of 2 to 4 m. The reason for this was that the specimens should have a slenderness factor of 2, i.e., the ratio between the length and diameter must be 2. This value was defined for a better comparison between the values of seismic wave speed and gain in UCS. Since the top two meters of the cores were more cracked, the crack-free length was not 200 mm. The second reasons was to perform a correlation between the P-wave velocity and uniaxial compressive strength (UCS). The compressive strength of the samples varied between 1800 to 4400 kPa, see Fig. \ref{Fig:17} which shows a good correlation between the compressive strength and the P-wave velocity (over 80\%), that is, the model explains 80\% of the variation in the data set. 

\begin{figure}[H]
 \centering
    \includegraphics[width=0.50\textwidth]{Fig_17.jpg}\\
  \caption{P-wave (compressiont) velocity as a function of UCS}\label{Fig:17}
\end{figure}

The compressive strength of the drill cores increases only insignificantly with depth, so the higher ambient pressure does not affect much the hardening process, as showed in correlation between the UCS and P-wave velocity in Fig. \ref{Fig:18}. 

\begin{figure}[H]
 \centering
    \includegraphics[width=0.50\textwidth]{Fig_18.jpg}\\
  \caption{Compressive strength as a function of depth}\label{Fig:18}
\end{figure}

Thus, the data from the drill cores showed only a weak tendency in increase of the compressive strength with depth. This tendency might be very weak because it is only based on a sampling above the level of 2 m below the ground surface where samples did not show any correlation between the depth and strength, Fig. \ref{Fig:18}. The analysis of the variation in dry density of soil at different depths did not show any differences. In all cases, the variation in density was within $100 kg/m^3$, Fig. \ref{Fig:19}.

\begin{figure}[H]
 \centering
    \includegraphics[width=0.50\textwidth]{Fig_19.jpg}\\
  \caption{Density as a function of depth}\label{Fig:19}
\end{figure}

In this test, the moraine clay specimens of core collected from the three different sites of ground surface down to the bottom of the stabilized layer (to 3.15 m) were evaluated. The soil specimens were tested to assess the correlations between the UCS with values of P-wave velocity using free-free resonant method obtained from relevant experiments. The data was evaluated statistically, visualized and plotted to interpret the degree of correlation between the UCS and velocity of seismic P-waves, and the variability in their values.

\subsection{X-ray diffraction test}

The analytical method based on X-Ray diffraction (XRD) test was performed to quantify the mineralogical composition of specimens. Theoretically, the mineral content in a crystalline sample can be determined using X-ray diffraction both qualitatively \cite{NABIH2020109384,WANG2021105354,MURER2021m0800747,PARANTHAMAN2021107992} and quantitatively \cite{BUTLER2021104662,WANG2022,BUTLER2020114474}. In case of the quantitative case, the XRD is applied on the reference minerals which are known to be commonly present in a given type of soil \cite{NOWAK2018133}. In our case, none of the minerals has had a known, fixed proportion which could be a function of the internal standard. The synchronous light test was performed for the qualitative analysis of the mineral soil composition. Sampling and analysis were performed both for natural wet and dried samples, Fig. \ref{Fig:20} and Fig. \ref{Fig:21}. 

\begin{figure}[H]
 \centering
    \includegraphics[width=0.50\textwidth]{Fig_20.jpg}\\
  \caption{Dried sample A. Minerals: chlorite (C), ettringite (E), feldspar (F), gypsum (G), illite (I), kaolinite (K), metavivianite (M), smectite (S), zeolite stilbite (Z)}\label{Fig:20}
\end{figure}

The measurements were carried out at the diffraction beam tube BLI 711 at MAX-lab, Öle Römers väg, Lund. The samples were measured at high resolution of 0.997Å wavelength from $4-40\degree 2 \theta$ angle. The measurements enabled to analysis the qualitative composition of soil. The evaluation of the contain of the samples was based on the qualitative X-ray analysis, i.e., without correct ratios. From the drill cores, samples were taken as one from a superficial layer and one from a deep layer. In total, four samples were examined. The most important minerals in these samples were identified via the search/match method. Using this method, the following minerals have been identified: chlorite (C), ettringite (E), feldspar (F), gypsum (G), illite (I), kaolinite (K), metavivianite (M), smectite (S), zeolite stilbite (Z), and amphibole (A)m, Fig. \ref{Fig:20} and Fig. \ref{Fig:21}. 

The purpose of the X-ray diffraction analysis was to identify the possible differences in the material composition of soil specimens collected in different boreholes. The samples were tested both for those with natural water content and for the oven-dried material. The pozzolanic reaction materials in the tested materials were not identified. This is probably due to the too small levels in relation to the high levels of clay minerals. Moreover, there were no differences between the top and bottom layers which could be seen in the tested soil. This proves that soil was homogeneous over the tested area.

\begin{figure}[H]
 \centering
    \includegraphics[width=0.50\textwidth]{Fig_21.jpg}\\
  \caption{Native sample B. Minerals: amphibole (A), chlorite (C), feldspar (F), gypsum (G), illite (I), kaolinite (K), metavivianite (M)}\label{Fig:21}
\end{figure}

\section{Discussion}\label{sec:6}

As demonstrated in this paper, the determination of deformation and compressive strength is possible using the estimated velocities of the elastic waves. Our study shown the advantages of the integrated use of the engineering geological techniques of borehole core sampling and geophysical seismic methods of data post-processing for geotechnical studies in industrial engineering. Measuring uniaxial compressive strength and stiffness of the compacted soil is necessary to ensure that it meets the requirements of the high bearing capacity and acceptable stiffness to the acting external loads during the exploitation of the building. A sophisticated method of X-Ray diffraction (XRD) based on the analytical tests was performed to quantify the mineralogical composition of specimens, as in existing similar studies \cite{Pires,Naik}. The soil analysis was exerted to alleviate the selection of binders and remove possible negative impact of the irrelevant binders and thus to optimise the workflow process of soil stabilization.

The results show an advanced and well-functioning approach combining geophysical methods of drilling, borehole drilling, acoustic measurements, seismic tests by P- and S-waves, soil analysis and stabilization. The presented results meets the design requirements regarding soil quality intended for foundation construction. Stabilization by GGBFS and CaCO has proven to be well suited for solidification of the fine-grained homogeneous soil. The stabilization enabled to significantly reduce the concrete slab in the construction polygon with the same or better performance. In connection with the construction works, a continuous quality control of each stabilized layer was carried out to monitor the quality of soil. 

A notable difference between our proposed method of soil testing with significant usage of geotechnical tools and most of the existing soil evaluation traditional laboratory approaches \cite{Barwar,Sharma,Lindh202302} is that it is build upon integrated scheme of methodology which incorporates the fundamental geophysical methods of the acoustic tests to solve engineering tasks on empirical evaluation of soil strength for testing geotechnical parameters of foundation for the Max IV facility of Lund University. The integrated method of geological borehole drilling presented in this paper includes the geophysical measurements and acoustic tests using seismic waves used together in the workflow of the project. Such integration of methods has proven to perform well with regard to the soil treatment technology, construction industry, material technology and economics, as also reported in similar studies for road base application using combination of the approaches \cite{Gupta,Lindh202234,Eberemu,Lindh202303}. 

Thus, the geological, geophysical and geotechnical approaches methods that use different approaches to the target research object were combined, i.e. soil, yet overlap in background principles of data measurements and processing. In contrast to the existing studies on soil stabilization, these methods were linked in a connected workflow that facilitated complex soil investigation in the project area. For the more detailed analysis of soil properties, the synchronous light test by the XRD was applied, which enabled the qualitative analysis of mineral composition of soil. Specifically, the specimens were first mined using the drilling of core samples as cylindrical specimens and processed data. The properties of soil samples were estimated through the analysis of the borehole logging and detected certain differences in the crosshole measurements based on the three core samples. The soil samples were tested using the traditional approach of grain distribution analysis and water content. 

The method of XRD was adopted for mineral analysis of specimens based on the heterogeneous crystalline structure of soil. The UCS and seismic tests were integrated, which presented the two approaches, the traditional evaluation of soil, which is similar to the existing studies \cite{Yilmaz}, and the non-destructive precise measurements that are applied in engineering works on soil improvement \cite{Staab,Madhyannapu}, for cross-correlation of the UCS and the elastic wave speed. The evaluation of the results has been performed using digital data processing, computing, visualization and plotting for decision making regarding the soil suitability as a subsoil foundation able to safely resist the load from the load bearing structures during the construction of Max IV facility of Lund University. 

The results of this study showed that no significant differences between the superficial and deep-lying samples could be seen in the investigated matrix, despite some natural small variations in the basic material specimens. The demonstrated results from the performed works on this project show great possibilities to stabilize a fine-grained soil. Specifically, this work tested the stabilization of the clay moraine used as a basis for constructions  with high requirements on sensitivity to vibrations, as well as high strength and stiffness.

\section{Conclusion}\label{sec:7}

Evaluating the soil properties during the subsoil construction cycle is essential for assessment of soil quality and compliance with the technical requirements. This includes the evaluation of parameters of soil that indicate whether the proposed substructure or building foundation corresponds to the technical requirements of stability and safety in terms of bearing capacity and stiffness. The extensive experimental tests using the \emph{in-situ} data collected from the boreholes were conducted with promising results of application of the proposed workflow for similar projects designed for the facilities with the advanced requirements for soil ground. 

This paper demonstrated the geotechnical exploration of soil with examples of application of geophysical methods. 

Future applications of the presented study can adopt presented methods of geotechnical exploration of the soil ground intended for construction of complex facilities. Similar works may also consider the described methods which were performed in fieldwork and laboratory. It is recommended to evaluate the results of the UCS using velocities of elastic P-waves, since these methods demonstrated to be effective and robust for control quality of strength of the stabilized soil. The integrated approach demonstrated in this study can be applied in further works using the benefits of the integrated seismic and geophysical methods in soil exploration studies. Specifically, such methods shown to be useful in civil engineering for testing quality of soil materials.

The presented results validate the following conclusions: 1) it is beneficial to apply geophysical methodology and tools for civil engineering works due to the demonstrated robustness and precision of data processing cycle: data capture, data analysis, data evaluation, data visualization and plotting, data interpretation and decision making. 2) the capability of data interpretation based on the analysis of the multi-format and heterogeneous data received from geophysical and geological exploration sources is realistic and advantageous for civil engineering and real geotechnical projects.

\vspace{10pt}

\noindent
\textbf{Author contributions}: All authors have accepted responsibility for the entire content of this manuscript and approved its submission.

\vspace{10pt}

\noindent
\textbf{Conflict of interest}: The authors state no conflict of interest.

\begin{acknowledgement}
We thank the anonymous reviewers for their valuable comments. The Division of Building Materials of Lund University provided the concrete cutter for sawing samples. The authors thank Ola Malmgren (Peab Anläggning AB) for project management and coordination, Brian Norsk Jensen, Dörthe Haase (Max IV Laboratory) for assistance with tests, Staffan Hansen (Polymer Technology, LTH) and Nils Rydén (Peab Grundteknik/LTH) for assistance with data interpretation. Roger Wisén (Ramböll), for assistance with interpretation of geophysical logging, Per-Olof Rosenkvist (LTH) for the assistance with the UCS tests. We also thank the reference group consisting of Richard Nilsson (Skanska Sweden AB), Roger Wisén (Ramböll), Brian Norsk (Jensen MAX IV Laboratory), Johan Blumfalk (Hercules Foundation). The project was financially supported by the Development Fund of the Swedish Construction Industry (Svenska Byggbranschens Utvecklingsfond, SBUF) as project 13152, the Nordic Community Builder Peab, Swedish Geotechnical Institute (SGI) and Max IV Laboratory of Lund University, Sweden. 
\end{acknowledgement}

%\newpage
%\bibliographystyle{IEEEtran}
\addcontentsline{toc}{section}{References}
\bibliographystyle{plain}
\bibliography{References_JMBM.bib}
%\nocite*{}

\end{document}